\label{sec_dpfs}
\subsubsection{Point Functions}
On a high level view, a point function $f$ is defined around a target index $\alpha$ and its target value $\beta$.
For each input $i\neq\alpha$ the point function will output $0$.
Only if the input $i$ is equal to the target index $\alpha$, $f$ will output $\beta$.
\[
f_{\alpha,\beta}(i) = \begin{cases} 
      \beta & \text{if } i = \alpha, \\
      0 & \text{else}.
   \end{cases}
\]
\cite{boyle2023information} describe a point function with $f_{\alpha, \beta}:[N]\rightarrow G$. 
$G$ is an Abelian group and $N$ is the domain size.
Note that $\alpha\in[N]$ and $\beta\in G$.
$f$ evaluates $\alpha$ to $\beta$ and to $0\in G$ for all other inputs.
Thus, $f_{\alpha,\beta}=(N,G,\alpha,\beta)$ represents the valid point function $f$.

In this research all following point functions will evaluate to $1$ for a given target index $\alpha$.
So:
\[
f_{\alpha,1}(i) = \begin{cases} 
      1 & \text{if } i = \alpha, \\
      0 & \text{else}.
   \end{cases}
\]
This implies that this special case $f_{\alpha,1}$ is a functional representation of a one-hot vector which has its only value unequal zero at index $\alpha$.
Such representations will be necessarry for representing the target index of which element to retreive from the database lateron.

\subsubsection{Distributed Point Functions}

\subsubsection{Boyle-Gilboa-Ishiai DPF Construction}